% page 441 of the facsimile (image p-440 of the scan)
% block 170.2 x 242.0 mm ; left margin 31.8 mm
\pgc{10.160mm}{0.000mm}{
\l{\cel{57}433}
\l{}
\l{}
\l{\sou{perihelio}. (\sou{astr.)}\cel{3}Punto ek la orbito di planeto ube olca esas}
\l{\cel{4}maxime proxima a Suno. - DEFIS.}
\l{}
\l{\sou{perikardio}. (\sou{anat.)}\cel{3}Sako membrana qua envelopas la kordio.-DEFIS}
\l{}
\l{\sou{perimetro}. Konturo qua limitizas spaco-parto determinita. - DEFIRS.}
\l{}
\l{\sou{perineo}. (\sou{anat.}) Regiono inter la organi sexuala e la anuso. -DEFIS}
\l{}
\l{\sou{periodo}. Tempo-quanto dum qua kozo parkuras la fazi di sua duro.}
\l{\cel{4}DEFIRS.}
\l{}
\l{\sou{periosto.}\cel{2}(\sou{anat.)}\cel{2}Membrano fibroza qua parkovras la osti.-DEFIS.}
\l{}
\l{\sou{periostito}. (\sou{patol.}) Inflamuro di la periosto. - DEFIS.}
\l{}
\l{\sou{peripatetiko}. La persno qua adheris a la filozofio di Aristoteles}
\l{\cel{4}- DEFIRS.}
\l{}
\l{\sou{peripecio}. Bruska iro-chanjo di la eventantajo ad altra. - DEFIRS.}
\l{}
\l{\sou{perisar}. (\sou{netrans.}) Supresesar per morto violentoza. EFIS.}
\l{}
\l{\sou{periskopo}. Optiko-tubo, uzata dat-unesme da la submarnavi, kom}
\l{\cel{4}aparato qua posibligas vidar dum submerseso. - DEFIS.}
\l{}
\l{\sou{peristalta}. (\sou{fiziol.})\cel{1}\sou{Movo peristalta}\cel{1}: Kontrakto, de supre ad}
\l{\cel{4}infre, di la fibri cirklatra di la digesto-tubo, qua havas}
\l{\cel{4}kom efekto pulsar kontinue adavan la nutrivi. - DEFIS}
\l{}
\l{\sou{peristilo}. (\sou{arkitekt.})\cel{2}Edifico di qua la periferio interna}
\l{\cel{4}kontenas kolonaro. - DEFIRS.}
\l{}
\l{\sou{peritoneo}. (\sou{anat.)}\cel{2}Membrano seroza qua esas quaze la tapeto}
\l{\cel{4}di la kavajo abdomina, ed envelopas la viceri quin ica kon}
\l{\cel{4}tenas. - EFIS.}
\l{}
\l{\sou{perjurar}. (\sou{netrans.}) I. Agar trompo-juro. - II. Violacar sua}
\l{\cel{4}juro-promiso. - EFIS.}
\l{}
\l{\sou{perko.}\cel{2}(\sou{zool.}) Fisho qua vivas en aquo sen-sala, e havas flosi}
\l{\cel{4}dornoza. - L.\cel{1}\sou{perca.}\cel{2}- EFIS.}
\l{}
\l{\sou{perkalo.}\cel{1}Kaliko dina, di qua la lanugo deprenesis. - DEFIRS}
\l{}
\l{\sou{perkalino.}\cel{2}Perkalo tintita e glata. - DEFIRS.}
\l{}
\l{\sou{perkutar.}\cel{1}(\sou{trans.)}I. (medicino). Explorar per la metodo qua kon}
\l{\cel{5}sistas ye frapar sur la parieto di korpo-kavajo por saveskar}
\l{\cel{5}(segun la naturo di la bruiso) quan esas la stando di la or-}
\l{\cel{5}gani en ta kavajo. - II. (\sou{armo})\cel{2}(\sou{aludante fusilo}). Agar la}
\l{\cel{5}frapo qua - per explozo di la pulvero - deskuplos la projek-}
\l{\cel{5}tilo di la pafarmo. - DEFIS.}
}
